% =====================================================================
\chapter{The Falsifiable Research Programme}\label{ch:program}
% =====================================================================

Every claim in this paper is no more than an interpretive framework unless it is converted
into testable form. This chapter performs that conversion, following the distinction
between hard core and protective belt in the methodology of research programmes
\cite{lakatos1970}: claims that are not load-bearing belong in the protective belt, not
reinforced into the hard core.

\section{Methodology: From Interpretive Framework to Testable Hypotheses}

\begin{principle}[Demotion procedure]\label{pr:demote}
For every claim, carry out the following operations: write down its \emph{observable
consequences}; write down what observation would \emph{refute} it; write down the
\emph{minimal experiment} that would yield that observation; and write down the
\emph{control condition} of that experiment. Any claim for which these four steps cannot be
completed is to be marked explicitly as an interpretive framework rather than an empirical
assertion.
\end{principle}

By Principle \ref{pr:demote}, the claims of this paper fall into three classes:

\begin{itemize}
\item \term{Proved}: propositions \ref{prop:oblivious}, \ref{prop:overconstrain},
      \ref{prop:stopping}, \ref{prop:leakpair}, \ref{prop:soundness},
      \ref{prop:active}, \ref{prop:asym-struct}, \ref{prop:operad-compose},
      \ref{prop:transport}, \ref{prop:observability}, \ref{prop:exponent}; theorems
      \ref{th:integration}, \ref{th:six}, \ref{th:centralise},
      \ref{th:condorcet}; proposition \ref{th:correlated}; corollaries
      \ref{cor:threshold}, \ref{cor:buildadopt}, \ref{cor:diversity},
      \ref{cor:three-failures}.
\item \term{Argued under explicit assumptions}: principles \ref{th:compression},
      \ref{th:necessity}, \ref{th:separation}, \ref{th:design}; propositions
      \ref{prop:no-godel}, \ref{prop:three-grounds}.
\item \term{Advanced as falsifiable hypotheses}: the hypotheses $H_1$--$H_{14}$ below. In
      addition, the structural proposal of Remark \ref{rmk:fibration-status} (promoting the
      level family to a strict fibration) and the renormalisation analogy of Remark
      \ref{rmk:rg} are not counted among the theorems; their status is that of a proposal
      and a heuristic.
\end{itemize}

\section{Fourteen Falsifiable Hypotheses}

\begin{hypothesis}[Structural convergence above a random baseline]\label{hyp:h1}
After controlling for system size, domain, platform, programming language and framework,
large software systems from different domains exhibit structural similarity on their
abstract skeleton (Definition \ref{def:skeleton}) significantly above a random baseline.
\emph{Operationalisation}: construct a collection of architecture graphs and skeleton maps
by the five-step design of \S\ref{sec:baseline}, take the structural metrics of the third
step as the feature vector, and test whether the similarity distribution of cross-domain
pairs stochastically exceeds a null model, such as degree-sequence-preserving graph
rewiring.
\emph{Falsification condition}: the similarity distribution does not differ significantly
from the null model.
\emph{Risk}: freedom in the choice of skeleton map; the map and the metrics must be
pre-registered.
\end{hypothesis}

\begin{hypothesis}[Convergence strength monotone in domain binding]\label{hyp:h2}
Conjecture \ref{conj:strength} holds: structural similarity decreases monotonically in the
domain binding $\mathrm{Dom}$ (Definition \ref{def:dom}).
\emph{Operationalisation}: annotate $\mathrm{Dom}$ at each level (estimated from the
fraction of domain axioms on which that level's semantics depends) and regress similarity
on $\mathrm{Dom}$.
\emph{Falsification condition}: the monotonicity fails, or the coefficient is not
significant.
\end{hypothesis}

\begin{hypothesis}[Phase-transition threshold determined by $n \times m$]\label{hyp:h3}
Whether a system builds a unified substrate is predicted better by the product of the
number of live entities and the number of shared resources (together with the frequency of
recurrent problems and the frequency of evolution) than by lines of code.
\emph{Operationalisation}: for a set of open-source projects, annotate manually or
semi-automatically whether a self-built implementation of any of the six functions appears;
fit separate logistic regressions on a proxy for $n \times m$ (order of magnitude of
concurrent entities, number of kinds of shared resource) and on lines of code; compare
their explanatory power.
\emph{Falsification condition}: lines of code explain the outcome no worse than
$n \times m$.
\emph{Risk}: automatic estimation of $n$ and $m$ is difficult; the number of concurrent
entities and of mutexes, connection pools, caches and the like can serve as a first proxy.
\end{hypothesis}

\begin{hypothesis}[Read-set scaling outperforms traditional size metrics]\label{hyp:h4}
The change read-set $\mathbb{E}[\rho]$ (Definition \ref{def:readset}) and its \emph{scaling
exponent} $\alpha$ with respect to system size (Proposition \ref{prop:exponent}) correlate
with defect rate and change failure rate better than traditional size metrics such as
module count, lines of code and cyclomatic complexity, and better than the current magnitude
of the read-set $\mathbb{E}[\rho]$ itself.
\emph{Operationalisation}: on several open-source projects with a complete commit history
and defect tracking, estimate the magnitude and the scaling of the read-set over the same
time windows, regress each in turn on defect density in a panel design, and compare their
explanatory power.
\emph{Falsification condition}: the scaling provides no predictive power beyond the
magnitude and the existing metrics.
\end{hypothesis}

\begin{hypothesis}[Layered model plus active probing outperforms exhaustive reading]\label{hyp:h5}
Under a constraint on unit cost (token budget or person-hours), a comprehension process
that uses a layered cognitive model plus active probing attains a higher task success rate
(scored against (U1)--(U5) of Definition \ref{def:understand}) than exhaustive reading.
\emph{Operationalisation}: given the same unfamiliar code base and the same set of tasks
(predict the impact surface, locate a defect, add a feature), compare three strategies:
sequential reading, layered summaries only, and layered summaries plus a probing loop.
Measure task success rate and success rate per unit cost.
\emph{Falsification condition}: the active-probing strategy is not significantly better, or
its advantage comes entirely from the number of probes rather than from the structure of the
strategy.
\end{hypothesis}

\begin{hypothesis}[Leak pairs are automatically discoverable]\label{hyp:h6}
The leak pairs of Proposition \ref{prop:leakpair} can be found by automated differential
testing, and the discovery rate is negatively correlated with module boundary quality.
\emph{Operationalisation}: for several libraries, generate execution pairs that are
``identical at the interface and different in internal state'', detect divergence in
external behaviour, and compare against human-annotated boundary quality scores.
\emph{Falsification condition}: automatically generated execution pairs fail to trigger
leaks effectively, or the discovery rate is unrelated to boundary quality.
\end{hypothesis}

\begin{hypothesis}[Specification quality is the present bottleneck]\label{hyp:h7}
At the present state of the art, the principal bottleneck on the path from intent to
verifiable execution is specification quality, not code generation capability or proof
search capability.
\emph{Operationalisation}: take existing benchmark evidence (specification generation
tasks, evaluations that separate code correctness from specification quality, proof success
rates) as initial evidence, and on a new controlled benchmark compare the marginal
contribution of each of the three links to end-to-end success rate.
\emph{Initial evidence}: a component-wise evaluation on 189 Lean tasks reports, for the
best model, code correctness 72.6\%, specification soundness and completeness 52.3\%, and
proof success rate 4.9\% (single trial) \cite{verina2026}; an evaluation of 245
operating-system kernel specification generation tasks shows that models remain deficient
in long context and specification completeness \cite{osvbench2026}; an evaluation that
instantiates specifications as trusted tests points out that proof difficulty interferes
with the assessment of specification quality \cite{coins2026}; a traceable refinement
method reports accuracy gains on the specification-generation subtask
\cite{verispecgen2026}; a mutation-driven method reports precision and recall
\cite{specsyn2026}.
These are empirical results on specific benchmarks, not general reliability theorems, but
they support the judgement that the three dimensions should be evaluated separately.
\emph{Falsification condition}: with the capability of the other links controlled, the
marginal contribution of specification quality is not significantly greater than that of
the other links.
\end{hypothesis}

\begin{hypothesis}[Consensus reliability is dominated by correlated error modes]\label{hyp:h8}
On intent-attestation tasks, the reliability gain from adding verifiers is bounded above by
the correlated error mode $q$ (Proposition \ref{th:correlated}); introducing heterogeneous
verifiers (people of different backgrounds, models from different sources, adversarial
roles) improves reliability more than merely adding homogeneous verifiers.
\emph{Operationalisation}: construct intent-attestation tasks, configure several verifier
groups (homogeneous or heterogeneous), measure how the error rate of the majority judgement
varies with group size and heterogeneity, and compare with the predictions of Equations
\eqref{eq:chernoff} and \eqref{eq:correlated}.
\emph{Falsification condition}: the rate at which the error rate of homogeneous verifiers
decreases with $n$ does not differ significantly from that of heterogeneous verifiers (that
is, $q$ is in fact small).
\end{hypothesis}

\begin{hypothesis}[The artefact binding gap is a measurable risk]\label{hyp:h9}
The binding gap (Equation \eqref{eq:binding}) between the verified artefact and the
artefact actually executed is a measurable risk in real deployments and can be
significantly reduced, and reducing it yields an observable reliability benefit.
\emph{Operationalisation}: introduce content addressing, reproducible builds and
deploy-time verification into several deployment pipelines, and measure the change in the
proportion of build nondeterminism, the rate of deployment drift events and the rate of
related incidents.
\emph{Falsification condition}: the gap is not measurable in practice, or the cost of
closing it far exceeds the benefit.
\end{hypothesis}

\begin{hypothesis}[The six functions are the surviving set under coarse-graining]\label{hyp:h10}
Under the coarse-graining sequence of Definition \ref{def:coarsegrain}, the structures that
survive are exactly the slots of $\mathcal{K}_6$ and their relational graph (Conjecture
\ref{conj:relevant}).
\emph{Operationalisation}: apply mechanical coarse-graining to several repositories
following Principle \ref{pr:coarsegrain-impl}, and record how the surviving set changes
with granularity.
\emph{Falsification condition}: the surviving set contains structures outside the six
functions, or some of the six functions are washed out already at intermediate granularity.
\emph{Risk}: the choice of coarse-graining threshold introduces an arbitrary parameter; the
threshold should be subjected to a sensitivity analysis.
\end{hypothesis}

\begin{hypothesis}[Incentive structure reduces correlated error]\label{hyp:h11}
On intent-attestation tasks, adopting an appropriate scoring rule (Principle
\ref{pr:scoring}) reduces the error rate of the majority judgement more than merely
increasing the number of verifiers, and a measurable fall in error correlation accompanies
it.
\emph{Operationalisation}: on the same set of specification-attestation tasks, compare
three configurations: few unincentivised verifiers, many unincentivised verifiers, and few
verifiers under a scoring rule; measure the error rate and the correlation coefficient of
errors between verifiers.
\emph{Falsification condition}: neither the error rate nor the error correlation
coefficient of the scoring-rule configuration is better than that of the
many-unincentivised configuration.
\end{hypothesis}

\begin{hypothesis}[Rate--distortion slope predicts refactoring benefit]\label{hyp:h12}
The slope $\partial L_T / \partial \DL$ of the rate--distortion curve (Remark
\ref{rmk:rd-refactor}) predicts which refactorings reduce defect rate or rework rate in
practice.
\emph{Operationalisation}: for a set of completed refactorings, estimate ex ante the
direction of displacement along the curve and the slope, and ex post regress the changes in
defect rate and change failure rate on them.
\emph{Falsification condition}: the slope provides no predictive power beyond naive
variables such as the size of the change.
\end{hypothesis}

\begin{hypothesis}[The trajectory of proof cost per unit of code]\label{hyp:h13}
Chapter \ref{ch:objections} declines to claim that formal verification becomes cheap automatically
and asks instead for a measurable cost metric. The hypothesis is that the ratio
\[
\kappa_{\text{proof}} \;=\; \frac{\text{effort to establish a stated guarantee}}{\text{size of the code carrying it}}
\]
is on a measurable trajectory, and that the trajectory---not its current value---decides whether
mechanised verification becomes routine engineering or remains confined to safety-critical work.
\emph{Initial evidence}: the seL4 functional-correctness proof is reported at about eleven
person-years for about 8{,}500 lines of C, or roughly $1.3$ person-years per thousand lines,
\emph{excluding} investment in reusable tools and frameworks \cite{klein2009}; figures that include
that investment are several times larger. The two figures differ by a factor worth stating, because
the reusable-tools share is precisely the part that should fall if the hypothesis holds.
\emph{Operationalisation}: for a fixed class of properties and a fixed prover, record
$\kappa_{\text{proof}}$ per project and per year, separating the reusable-infrastructure share from
the per-project share; fit the trajectory over the observation window.
\emph{Falsification condition}: no measurable trend, or a trend attributable entirely to the
reusable-infrastructure share while the per-project share is flat. Since the paper declines to
predict which way this goes, both outcomes are informative.
\emph{Risk}: the denominator is contentious---lines of code, of specification, or of proof---and the
choice must be pre-registered or the trajectory is not comparable across projects.
\emph{Interpretation bands}: the thresholds conventionally proposed for reading the trajectory, such
as $0.1$ for ``routine'' and $0.5$--$1$ for ``specialist'', are stipulated decision bands rather
than measured quantities, and are labelled as such here.
\end{hypothesis}

\begin{hypothesis}[Agentic generation can lower the read-set scaling exponent]\label{hyp:h14}
Proposition \ref{prop:exponent} makes the read-set scaling exponent $\alpha$ the quantity that
decides whether a system's governance cost grows linearly or quadratically in its size. Every
hypothesis above concerns systems built by people or by human-directed tools. This one asks the
AI-era question directly: can an agentic generator, given licence to restructure, \emph{discover}
architectures with a lower $\alpha$ than the human baseline, or does it only move constants within
an architecture it inherits? The distinction is not cosmetic---it is the difference between
automation that accelerates work and automation that changes what remains possible at scale.
\emph{Operationalisation}: start an agent and a human control group from the same code base with the
same feature backlog and the same freedom to refactor; track $\mathbb{E}[\rho]$ and the fitted
$\alpha$ against system size over the same window; compare trajectories.
\emph{Falsification condition}: $\alpha$ does not differ significantly between arms, or the apparent
difference is explained entirely by the volume of code written rather than by its structure.
\emph{Risk}: $\alpha$ is estimated through the read-set proxy of Definition \ref{def:readset}, whose
validity is itself hypothesis $H_4$; a failure of $H_4$'s proxy would make this hypothesis
unfalsifiable rather than false, and the two must not be confused.
\end{hypothesis}

\begin{table}[htbp]
\centering
\footnotesize
\caption{Summary of hypotheses and experimental designs}\label{tab:hypotheses}
\begin{tabular}{L{0.9cm}L{4.0cm}L{4.6cm}L{4.6cm}}
\toprule
\textbf{ID} & \textbf{Hypothesis} & \textbf{Principal measurement} & \textbf{Principal risk} \\
\midrule
$H_1$ & Structural convergence above a random baseline & Skeleton similarity vs.\ null
  model & Excessive freedom in the choice of skeleton map \\
$H_2$ & Convergence strength monotone in domain binding &
  Regression of similarity on $\mathrm{Dom}$ &
  $\mathrm{Dom}$ is hard to annotate independently \\
$H_3$ & Threshold determined by $n\times m$ & Comparison of logistic-regression explanatory
  power & Automatic estimation of $n$ and $m$ \\
$H_4$ & Read-set scaling outperforms traditional size metrics &
  Panel regression of read-set scaling on defect density & Accuracy of the read-set proxy \\
$H_5$ & Layered model plus active probing outperforms exhaustive reading &
  Task success rate per unit cost & Subjectivity of tasks and scoring \\
$H_6$ & Leak pairs are automatically discoverable & Differential-testing discovery rate &
  Difficulty of triggering valid execution pairs \\
$H_7$ & Specification quality is the bottleneck & Marginal contribution of the three links &
  Ground truth and coverage of the benchmark \\
$H_8$ & Consensus dominated by $q$ & Error rates of homogeneous and heterogeneous verifiers &
  Construction and measurement of heterogeneity \\
$H_9$ & Artefact binding gap is measurable & Build nondeterminism and drift event rate &
  Availability of real deployments \\
$H_{10}$ & The six functions are the coarse-graining surviving set &
  Surviving set under the coarse-graining sequence &
  Arbitrariness of the coarse-graining threshold \\
$H_{11}$ & Incentive structure reduces correlated error &
  Error rate and error correlation coefficient &
  Experimental ethics and cost of incentive schemes \\
$H_{12}$ & Rate--distortion slope predicts refactoring benefit &
  Regression of defect rate before and after refactoring & Estimability of the slope \\
$H_{13}$ & Proof cost per unit of code follows a measurable trajectory &
  Fitted trajectory of $\kappa_{\text{proof}}$, tools share separated &
  Contentious denominator; comparability across projects \\
$H_{14}$ & Agentic generation can lower the read-set exponent $\alpha$ &
  Trajectory of fitted $\alpha$ against system size, two arms &
  Depends on the read-set proxy of $H_4$ \\
\bottomrule
\end{tabular}
\end{table}

\section{Three Priority Research Directions}

The three directions below connect requirements engineering with formal methods,
compilation correctness with resource analysis, and program contracts with runtime
security. They share one characteristic: each has a precisely stated question, a
falsifiable hypothesis, a method of testing and a principal risk.

\subsection*{Direction One: Specification-Quality Evidence Chain for Real Intent}

\paragraph{The precise question.}
Given a set $R$ of natural-language requirements, a domain model $D$ and a generated
specification $S$, how is an independent evaluation process to be built that checks
separately: requirement coverage; exclusion of forbidden behaviour; retention of permitted
behaviour; explicitness of environmental assumptions; and the satisfiability and
implementability of the specification?

\paragraph{Existing foundations.}
Requirements traceability research, requirement formalisation and verification, vacuity
detection, and recent specification-generation and joint-evaluation benchmarks. The key
connection is \emph{semantic traceability} in requirements engineering and
\emph{specification-versus-program counterexamples} in formal methods.

\paragraph{Falsifiable hypothesis.}
Compared with a single ``proof passed'' signal or a single ``similarity to a reference
specification'' score, atomic requirement traceability plus positive and negative testing
plus specification-strength and mutation checks plus independent domain review identifies
the \emph{under-constraint}, \emph{over-constraint} and \emph{intent omission} of a
specification earlier, that is, it is the concretisation of Hypothesis $H_7$.

\paragraph{Method of testing.}
On small benchmarks such as payments, permission control and idempotent interfaces,
establish requirement atoms annotated by multiple stakeholders, permitted and forbidden
trace sets, and environmental preconditions; compare four approaches: pure generation,
requirement decomposition, traceability testing, and mutation refinement with expert
assistance. Measure precision and recall, error categories, the rate of missed
prohibitions, the over-constraint rate, expert agreement and cost.

\paragraph{Falsification condition.}
If adding the traceability and mutation layers does not significantly reduce the rate at
which hidden negative requirements are missed, or if the gain comes only from leaked
reference test answers, the hypothesis fails.

\paragraph{Principal risk.}
The ``gold standard'' for real intent may itself be contested.

\subsection*{Direction Two: Joint Functional--Resource Translation Certificates for Actual Artefacts}

\paragraph{The precise question.}
Given a source program $P$, a functional contract $\varphi$, a resource budget $B$ and a
fixed target platform model $H$, can a certificate $C$ be generated for \emph{each
compilation} such that an independent verifier checks semantic preservation and analysable
resource constraints at the same time?

\paragraph{Existing foundations.}
Semantic-preservation proofs for verified compilers, per-translation validation,
translation certification, proof-carrying code, automatic amortised resource analysis and
worst-case execution time analysis. The substantive connection is to place the
``behaviour relation'' and the ``cost relation'', currently verified \emph{separately},
inside a \emph{cost-carrying behaviour semantics}.

\paragraph{Working model.}
The program semantics is extended to a pair $(\tau, c)$, where $\tau$ is an observable
trace and $c$ a cost; the translation proof should guarantee that target traces are
explicable by the source semantics, and should prove $c_{\text{target}} \le B$ under an
explicit cost model, or else give a conservative relation relative to the source cost.

\begin{remark}[An assumption that must be avoided]
It must not be presupposed that optimised compilation always lowers cost monotonically.
Optimisation may improve average performance while worsening the worst case; the cost
relation must therefore be a conservative relation comparable in both directions, not an
assumption of monotonicity.
\end{remark}

\paragraph{Method of testing.}
Choose a reproducible sublanguage and target architecture, and compare three routes:
verified compilation, translation validation, and ordinary compilation with post-hoc
verification. Record certificate size, verification time, coverage of program categories,
tightness of resource bounds, false-rejection rate, and robustness to changes in compiler
optimisation.

\paragraph{Falsification condition.}
If the resource proofs do not remain sound once calibrated against a real platform, or if
certificate generation and verification cost far more than the baseline with no safety
benefit, the scope should be narrowed to specific safety-critical domains.

\paragraph{Principal risk.}
The target processor microarchitecture and scheduling model are not precise enough.

\subsection*{Direction Three: Composition Contracts and Runtime Assurance for Untrusted Components}

\paragraph{The precise question.}
When a service is composed of verified components, generated components that are not fully
verified, and external services, how are assumption--guarantee contracts, information-flow
policies, resource budgets and runtime monitoring to be used to give an auditable overall
safety proposition together with an upper bound on failure impact?

\paragraph{Existing foundations.}
Modular assumption--guarantee reasoning, contract composition for dynamical systems,
information-flow non-interference, verified kernels, and studies of uncertainty in runtime
verification. The substantive connection is to place the formal contracts of subcomponents
and the runtime detection and enforcement of contract premises and guarantees within a
single system argument.

\paragraph{Falsifiable hypothesis.}
When every component has explicit assumptions, entry and exit events are observable, and
contract violations can trigger isolation or rollback from an independent monitor, a
composable conditional guarantee can be given for a specified safety property without
proving the internal implementation of each component correct.

\paragraph{Method of testing.}
Build a multi-component payment flow: an idempotent request service, an identity and
permission service, a generative risk-decision component, an external database and a
network simulator. Under fault injection, compare three groups (no contracts, static
contracts, contracts plus monitoring and rollback) and measure duplicate charges,
cross-user data leakage, detection latency, isolation success rate and performance cost.

\paragraph{Falsification condition.}
If the events observable by the monitor are insufficient to decide the critical properties,
or if the rollback mechanism cannot itself remain safe, the architecture cannot provide the
guarantee it claims.

\paragraph{Principal risk.}
For subtle semantic properties that go unobserved, runtime assurance can constrain only the
detectable and enforceable part.

\section{Six-Stage Implementation Roadmap}

\begin{table}[htbp]
\centering
\small
\caption{Research implementation roadmap}\label{tab:roadmap}
\begin{tabular}{L{1.0cm}L{6.3cm}L{7.0cm}}
\toprule
\textbf{Stage} & \textbf{Core tasks} & \textbf{Deliverables} \\
\midrule
1 & Define the benchmark and the threat model: choose a payment service and an idempotent
  interface as the running case; fix the assets, roles, observable events, forbidden
  behaviour, environmental assumptions and time limits &
  Set of requirement atoms, positive and negative trace sets, list of environmental
  assumptions (attested by stakeholders) \\
2 & Specification quality assessment: compare pure generation, traceability refinement, and
  mutation- and counterexample-driven approaches; check for under-constrained
  specifications, over-constrained specifications, vacuous satisfaction, unsatisfiability
  and unimplementability &
  Specification quality assessor; for each requirement identifier, the corresponding
  specification clause, test and attestation record \\
3 & Program synthesis in a restricted language: choose an explicit sublanguage and a logic
  of controlled decidability; separate candidate generation from proof checking &
  Candidate programs, proof obligations, counterexample classification; distinguish
  unsatisfiable, unknown, timeout and inexpressible \\
4 & Artefact and resource proofs: choose a mature compilation or translation-validation
  route; specify reproducible experimental hardware or a simulation model &
  Compilation-chain metadata, semantic-preservation or translation certificates, resource
  bounds and model boundaries \\
5 & Composition and runtime assurance: write assumption--guarantee contracts for external
  services and not-fully-verified modules; add independent monitoring, sandboxing, auditing
  and safe rollback &
  System-level contract composition results, fault-injection report, residual risk list \\
6 & End-to-end evaluation: measure assurance coverage, false negatives, false positives,
  development cost, proof cost, artefact verification cost, resource budgets and failure
  recovery &
  A public benchmark capable of falsifying the research hypotheses, a tool prototype and an
  experimental report \\
\bottomrule
\end{tabular}
\end{table}

\section{Success Criteria}

The criteria below assess whether the roadmap succeeds. They are designed to be
independently checkable:

\begin{enumerate}[label=(S\arabic*)]
\item \term{Explicit object}: the object of each guarantee (which artefact, which
      semantics, which environment) can be uniquely identified.
\item \term{Real assumptions}: the environmental assumptions relied upon can be tested, not
      assumed by default.
\item \term{Independent evidence}: the judgement of evidence does not depend on the
      generator's self-assessment.
\item \term{Distinguishable failure categories}: unsatisfiable, unknown, timeout,
      inexpressible, and candidate refuted by counterexample must be reported separately and
      not merged into a single ``failure''.
\item \term{Measurable cost}: generation cost, proof cost, verification cost and runtime
      cost can be measured separately.
\item \term{Falsifiability}: the refutation of at least one $H_i$ would weaken the core
      claims of the programme.
\end{enumerate}

\begin{remark}[Self-examination of the programme]
One core claim of this paper is that ``generation is untrusted and verification must be
independent''. The principle applies to the paper itself: its load-bearing propositions
must be independently checkable and must not depend on any unproved conjecture. Remark
\ref{rmk:q1q8} and the list of non-claims in Chapter \ref{ch:objections} are the concrete
results of that self-examination.
\end{remark}

\section{Minimum Viable Subset and the Adoption Risk}\label{sec:mvs}

The hypotheses above share a failure mode that none of them covers. They are each refuted or
confirmed by experiment, and the framework can also simply fail to be adopted, in which case no
experiment is run at all. The risk is worth stating explicitly because it is orthogonal to
falsifiability, and because it has a specific shape: the formal objects cost something to
maintain---declarations must be written, kept current and defended---and if that cost exceeds the
benefit, the practice dies of neglect rather than of refutation.

\begin{principle}[Minimum viable subset]\label{pr:mvs}
If only three elements of the framework can be adopted, they should be:
\begin{enumerate}[label=(\alph*)]
\item \term{the oblivious-set declaration}: each level or module states what it is permitted never to
      care about (Definition \ref{def:oblivious});
\item \term{the leak-pair test}: for each claimed boundary, a differential test seeking two executions
      with equal interface traces and different external behaviour (Proposition \ref{prop:leakpair},
      Corollary \ref{cor:diagnose});
\item \term{the level attribution of verification obligations}: each property is assigned to the
      lowest level that can express it, which decides whether it is mechanically decidable,
      monitorable, or a matter for attestation (Definition \ref{def:obligation-assign}).
\end{enumerate}
These three form the smallest closed loop in the framework: a declaration, a test of the declaration,
and a routing decision that consumes both. The contract operad of \S\ref{sec:operad} and the consensus
bound of \S\ref{sec:consensus} are refinement layers---valuable where the corresponding difficulty is
present, dispensable where it is not. Principles \ref{pr:descent} and \ref{pr:two-criteria} are what
make the loop usable: the first says how deep to look, the second says what the declaration has to
satisfy.
\end{principle}

The measurable form of the adoption claim is deliberately modest, and it is the metric Chapter
\ref{ch:levels} already proposes for its own empirical content: the rate at which properties that
should be assured mechanically are instead left to manual review, and the rate at which claimed
boundaries fail their own tests. If adopting the three elements above does not move those rates in a
real project, then the framework's problem is not that a hypothesis was false but that the remedy
costs more than the disease.

Stating the adoption risk as a fifteenth hypothesis would also be defensible, with adoption rate and
cost-to-benefit ratio as its measurements. This paper does not, for a reason that is itself part of
the framework's discipline: the falsification of a hypothesis should turn on an observation about the
world, whereas non-adoption turns on a decision by practitioners, and letting one word cover both
would blur the distinction between verification and validation that Chapter \ref{ch:intent} works to
keep sharp. The choice is recorded here rather than made silently.
